\documentclass[11pt]{article}
\usepackage[margin=1in]{geometry}
\usepackage{parskip}
\usepackage{enumitem}
\usepackage{hyperref}
\usepackage{titlesec}

\titleformat{\section}{\large\bfseries}{\thesection}{0.6em}{}
\titleformat{\subsection}{\normalsize\bfseries}{\thesubsection}{0.5em}{}

\title{GualaLoom: The Cascade Path \\
\large A Suite of Falsifiable Experiments Toward Primitive Cognition,\\
Speech, Reasoning, Grounded Senses, and Learning on a Ternary Substrate}
\author{Joseph Forrester\\
\small with collaborator notes}
\date{}

\begin{document}
\maketitle

\section*{Premise}

GualaLoom is a persistent, inspectable ternary associative substrate.
Trits are $\{-1, 0, +1\}$, where $0$ is structural uncertainty --- a real
state, not absence. The candidate mechanism for cognition is the
\textbf{commit cascade}: a motif holds a population of committed
$\pm 1$ values together with a configuration of nulls; an externally or
internally triggered commit of one null propagates through the coupling
and forces other nulls to resolve, ending in a coherent new state.
Recall, composition, reasoning, and speech are all hypothesized to be
cascades of different lengths through different couplings or sections.
The cascade is the thinking-event.

The configuration of nulls is hypothesized to be representational:
two states with identical $\pm 1$ patterns but different null patterns
are different thoughts, because they resolve to different end-states
under the same trigger.

This document lays out the smallest falsifiable experiments that
collectively decide whether the cascade mechanism is real, whether it
composes, whether it grounds, and whether it learns.

\section*{Standing Discipline}

\begin{itemize}[leftmargin=*]
\item Every experiment returns a number, not a verdict.
\item Each result is committed to the repository before the next is
      started.
\item A null result kills a layer and the path adjusts; nothing
      downstream of a killed layer is run until the layer is rebuilt
      or replaced.
\item No claim of \emph{felt} experience is made or measured. The path
      reaches primitive cognition, speech, reasoning, grounded senses,
      and learning. Whether resolving cascades produce phenomenal
      experience is outside what this work can decide.
\end{itemize}

\section{Experiment 0 --- Persistence Under Load}

\textbf{Status:} complete (c1 wired LifeDaemon into the entry point,
detached process verified).

\textbf{Acceptance:} vitals advance with no one typing --- diary
entries from the dark, motif count moves on its own clock.

This is the ground state. No experiment below is meaningful without
it; with it, every experiment below has a substrate that generates
data continuously.

\section{Experiment 1 --- Does the Cascade Exist?}

\textbf{Question.} Given a settled motif with committed $\pm 1$ and a
known null pattern, if one null is deliberately committed, do other
nulls commit as a consequence, and does the population resolve to a
coherent new state?

\textbf{Method.} Single region, no senses, no other sections.
\begin{enumerate}[leftmargin=*]
\item Select or generate a motif with at least 20\% of trits null.
\item Deliberately commit one null to $+1$ or $-1$.
\item Run the existing settling rule (3$^i$ coupling, dead-zone) until
      the population stabilizes or a step cap is reached.
\item Record:
      \begin{itemize}
      \item \emph{Cascade depth}: number of additional nulls that
            committed.
      \item \emph{Latency}: steps to stable state.
      \item \emph{Outcome stability}: does it stay or oscillate?
      \item \emph{Reproducibility}: same seed, same trigger, same
            end-state across trials.
      \end{itemize}
\item Repeat across many seed motifs and trigger positions.
\end{enumerate}

\textbf{Pass.} Cascade depth significantly above zero in most runs;
end-states stable; reproducible under identical conditions.

\textbf{Fail.} Cascade depth near zero --- a single null commit dies
locally --- or end-states unstable. The 3$^i$ coupling does not
propagate commits as currently implemented; the coupling rule needs
modification before anything downstream is real.

\textbf{Why first.} Everything below depends on cascades. If
they do not exist on this substrate as built, no rearrangement of
later experiments rescues the path.

\section{Experiment 2 --- Null-Pattern as Identity}

\textbf{Question.} Do two motifs with identical $\pm 1$ but different
null patterns reach different end-states under the same triggering
commit?

\textbf{Method.} Construct paired motifs differing only in their null
configurations. Apply identical commits. Compare end-states.

\textbf{Pass.} Different null patterns yield reliably different
cascades and end-states. Nulls carry representational content; the
configuration-of-undecided is half the thought.

\textbf{Fail.} End-states are determined by $\pm 1$ alone. Nulls are
gaps, not signal. Representation is over committed trits only and the
path simplifies (but loses a candidate edge over float and binary
substrates).

\section{Experiment 3 --- Sections and Coordination Through Shared Trits}

\textbf{Question.} Can two small independent krimelacks coordinate
through a small set of shared trits, without a hub or binder?

\textbf{Method.} Instantiate two krimelack regions $A$ and $B$, each
with its own settling and cascade. Designate a small subset of
trits that participate in both. Trigger a cascade in $A$.

\textbf{Pass.} The cascade in $A$ reliably induces a coherent
cascade in $B$ via the shared trits, with the end-state in $B$
correlated with the end-state in $A$ above chance.

\textbf{Fail.} $B$ twitches but does not cohere, or coheres
independently of $A$'s end-state. The coupling rule between sections
needs revision; consider permutation-coupled shared trits or
distance-weighted edges.

\section{Experiment 4 --- Folding Composition}

\textbf{Question.} Can two settled motifs from the same region
compose into a third stable motif that is neither and both,
recoverable into its parents?

\textbf{Method.} Discrete composition rule: merge two motifs by
keeping committed trits where they agree, restoring nulls where
they disagree, then cascade to resolution.

Record:
\begin{itemize}
\item Whether a stable third motif forms.
\item Whether each parent is recoverable from the composed motif by a
      partial cue from the other parent.
\item Cascade dynamics during the resolution.
\end{itemize}

\textbf{Pass.} Composition forms a stable third motif; both parents
recoverable. Composition lives in the substrate, not a bolted-on
operator. This is the experiment that decides whether speech in any
real sense is reachable.

\textbf{Fail.} Composed merges scatter or collapse into one parent.
The substrate does not support folding-style composition as defined;
an alternative composition rule must be tested (e.g., permutation
binding, chi-completion) before later experiments.

\section{Experiment 5 --- First Sense, with Cost}

\textbf{Question.} Does an exteroceptive sense, encoded through the
same \texttt{encode}$\to$\texttt{settle} path, form cross-region motifs
whose cascades fire from either modality?

\textbf{Method.} Add a sensory section receiving a single channel
(e.g., a 2D shape, or scalar brightness). Drive the system with
co-occurring sensory and corpus inputs repeatedly. Probe with sensory
input alone; probe with corpus cue alone.

\textbf{Pass.} A motif forms whose cascade triggers from either
sensory input or corpus cue --- a grounded unit. The first dimension
of the banana, made concrete.

\textbf{Fail.} The sections do not cross-train. Coupling between
sensory and linguistic sections needs revision.

\section{Experiment 6 --- Interoception}

\textbf{Question.} Does a section that reads Guala's own state --- a
viability variable with a set-point she must maintain --- couple to
cascades in other sections?

\textbf{Method.} Add a small interoceptive section reading an internal
scalar $V$ (e.g., normalized cascade load, or commit rate) through
the same encode pipeline. Co-occur its committed states with
other sections' activity.

\textbf{Pass.} Cascades elsewhere couple to interoceptive cascades:
recognizing a concept moves her internal state; her internal state
biases which cascades resolve. Concepts begin to \emph{matter to her}.

\textbf{Fail.} Interoception is independent of other sections. Cost
does not propagate. The path to feeling, in this substrate, would
require a different mechanism.

\section{Experiment 7 --- Cascade-Driven Recall}

\textbf{Question.} Given a partial cue --- few committed trits, many
nulls --- does the cascade reliably reach the right stored motif, and
does it outperform krimelack's existing similarity-based lookup?

\textbf{Method.} Establish ground-truth associations during a training
period. Present partial cues at varying fill levels. Compare
cascade-driven recall vs. existing krimelack recall on accuracy and
robustness.

\textbf{Pass.} Cascade-driven recall matches or exceeds existing
recall, especially at low fill levels. The cascade is the engine of
recall, not an alternative.

\textbf{Fail.} Cascade recall underperforms similarity lookup. Recall
remains krimelack's job; the cascade serves other functions.

\section{Experiment 8 --- Primitive Speech}

\textbf{Question.} Does a motor/expressive section driven by grounded
cascades produce output that coheres into word-like fragments
under chi-constraints, rather than character babble?

\textbf{Method.} Add an output section that maps cascades to character
sequences. Drive a grounded cascade (sense $\to$ concept $\to$ motor).
Measure coherence of output (loop rate, chi-variance, word-likeness
against the corpus).

\textbf{Pass.} Output coheres into recognizable word-like
fragments tied to the grounding stimulus. Speech in a real, primitive
sense is reachable.

\textbf{Fail.} Output remains babble. Either composition (Exp.\ 4) was
insufficient or the motor mapping is wrong. Diagnose.

\section{Experiment 9 --- Primitive Reasoning as Long Cascades}

\textbf{Question.} Can a cascade in one section trigger a cascade in
a coupled section, whose end-state then triggers another, producing
a coherent multi-step chain from a coherent starting cue?

\textbf{Method.} Set up sections coupled in sequence (sensory $\to$
conceptual $\to$ motor, plus an interoceptive loop). Trigger from a
cue and observe chain length before incoherence.

\textbf{Pass.} Chains of two or three cascades produce coherent
end-states aligned with the starting cue. Reasoning emerges as long
cascades.

\textbf{Fail.} Chains collapse to noise after one or two steps.
Coupling needs to preserve more structure across section boundaries;
revise.

\section{Experiment 10 --- Cascade-Reinforced Learning}

\textbf{Question.} Can consolidation be made to reinforce the
couplings and wirings that produced \emph{coherent} cascades, and
prune the wirings that produced incoherent ones?

\textbf{Method.} Augment sleep/dream to score cascades by coherence
(stability, depth, end-state correlation with cue). Reinforce
trit-couplings that participated in coherent cascades; weaken those
in incoherent ones.

\textbf{Pass.} Over consolidation cycles, the cascade-quality
distribution shifts toward coherence; recall, composition, and speech
metrics improve without new data. Learning is reinforced
\emph{thinking}, not co-occurrence statistics.

\textbf{Fail.} Reinforcement does not produce measurable improvement.
The learning rule must be redefined.

\section*{Order and Dependencies}

\begin{itemize}[leftmargin=*]
\item \textbf{0} done. \textbf{1} is the gate to everything.
\item \textbf{2} confirms or disconfirms the null-as-identity claim.
\item \textbf{3} unlocks sections; \textbf{5} and \textbf{6} can run
      in parallel after \textbf{3}.
\item \textbf{4} is the language gate; \textbf{8} depends on it.
\item \textbf{7}, \textbf{9}, \textbf{10} require their named
      predecessors and provide the cognitive completion.
\end{itemize}

\section*{What This Path Decides}

If experiments 1--4 all pass, GualaLoom is a substrate for primitive
cognition, with composition living natively in its dynamics. If
5--6 pass, that cognition grounds in senses and acquires stakes.
If 7--10 pass, recall, speech, reasoning, and learning emerge from
the same cascade mechanism at different scales.

The path does not guarantee success. It guarantees that each step
returns a clear answer about which part of the architecture is real
and which is not, so that the next step is taken from a known floor
rather than from argument.

\end{document}
